\section{LITERATURE REVIEW}

The development of robust, automated weed detection systems has been a central focus of precision agriculture research over the past decade. Making a universal weed detection framework is hindered by technical challenges, such as environmental variability, lack of annotated data, and the visual similarities between crops and weeds. This section critically examines the use of different computer vision techniques applied to this domain over the years. It progressed from early supervised machine learning approaches to the current state-of-the-art in artificial intelligence.

\subsection{Traditional Deep Learning in Precision Agriculture}
The shift from traditional machine vision to deep neural networks is one of the recent advances that increased the accuracy of automated weed recognition under controlled environments. Supervised Convolutional Neural Networks and one-stage object detectors, such as Yolo family networks, allow to automatically learn higher-level features from raw pixel data \cite{abedin2026} \cite{hamrani2025}. Recent developments in neural network design have shown outstanding performance in terms of speed and accuracy. According to \cite{sharma2024}, Yolo v11 achieved 13.5 ms inference speed and 89.10\% mAP on average. In turn, Yolov4-Tiny achieved 42.24 FPS on UAV maize weed detection \cite{pei2022}. The accuracy of a custom-made 5-layered Convolutional Neural Network reached 95.12\% with only 16.754 ms of inference speed and 0.012 GB memory consumption on a Raspberry Pi platform \cite{razfar2022}.

However, the accuracy of such methods is limited by the quality and quantity of annotated data, which is a known limitation of traditional computer vision (Deng et al., 2024). To train such a model, one needs to create pixel-level annotations, which is a very time-consuming process. According to (Gao et al., 2025), annotating one image could take 15-20 minutes. Therefore, it is not feasible to collect annotated data for every type of weed on the planet.

In addition to the lack of annotated data, deploying such a model in real-world conditions faces the problem of domain shift. In other words, if the distribution of data where the model is deployed differs from the distribution of data on which it was trained, its accuracy drops significantly (gao et al). The magnitude of this drop was recently quantified in an attempt to classify crops using a Source-Only approach: “On the cross-continental crop classification task, the mAP of a Source-Only model decreased by as much as 33.67\% from a value of 99.17\% (upper bound) to 65.50\% on the target set, completely failing to recognize the “Cotton” class” (Li et al., 2025). Similar results were presented when analyzing seasonal changes: “The baseline YOLOX, YOLOv8 and DINO detectors experienced performance drops ranging from 8.8\%, to 14.5\% on the cross-year experiment, with accuracy decreases of up to 26.1\% for DINO on the maize classification task” (Deng et al., 2024). Domain shift poses an even bigger challenge when it comes to recognizing plant species in the field. On images with a high degree of overlap and varying lighting conditions, “the mAP of Faster R-CNN fell to as low as 9.55\%” (Gao et al., 2025).


\subsection{Unsupervised Domain Adaptation (UDA) and Semi-Supervised Learning}
In order to save the astronomical amount of manual labor required to annotate a large number of images, and to reduce the catastrophic failures that occur due to domain shift, the agricultural research community has increasingly shifted to Unsupervised Domain Adaptation (UDA) and Semi-Supervised Learning (SSL) \cite{nadeem2026}. Initial methods based on Generative Adversarial Networks (GANs) and CycleGAN for style transfer produced results that were adequate for getting the source images translated to fit into target domains, with an Intersection over Union (IoU) increase of around 10\% on various source-target field pairs for the mean Intersection over Union (mIoU) metric \cite{ilyas2023} \cite{magistri2023}. To fill the significant visual gap between various plant images from the internet and the limited target domain agricultural robot data, Multi-level Attention-based Adversarial Discriminators (MAAD) have been incorporated into feature extractors, which resulted in 7.5\% higher object detection precision on unlabeled target domains \cite{tzouras2023}.

Likewise, semi-supervised teacher-student architectures have shown that using large amounts of unlabeled data can greatly improve accuracy. YOLOv8l-based WeedTeacher achieved a 2.6\% improvement over the fully supervised baseline with only 10\% of the labeled vegetable field (deng et al). However, other semi-supervised methods based on FCOS and Faster R-CNN on multi-class cotton datasets had about 76\% and 96\% detection accuracy, and greatly reduced the dependency on annotations (li et al). To avoid the model being overfitted by predicting false targets, class covariance analysis based pseudo labelling has been proposed \cite{huang2024}. New feature-level UDA improvements have enabled the development of masked adversarial alignment (CCUDA-SS) to selectively adapt models to uncertain areas while incurring minimal runtime complexity, resulting in up to +5.23\% mIoU improvement on eight datasets \cite{nadeem2026}. Moreover, 80.53\% mIoU was achieved by the 1-shot Supervised Domain Adaptation (SDA) model, which was very close to the 83.69\% accuracy of the fully supervised network when UDA was combined with the use of augmentation scheduling \cite{ilyas2023}.

But there are inherent limitations of UDA and SSL approaches, which limit their universal application. The most important weakness is "pseudolabel drift," which occurs when the domain changes significantly, and when a label is misclassified at the beginning, it is continually reinforced, actively contaminating the teacher model's target distribution \cite{nadeem2026}. Empirical results demonstrate that sub-optimal SSOD framework designs can negatively impact performance when domain shift is strong: for example, the DenseTeacher framework was applied to cross-domain weed datasets, yielding a significant decrease in mAP@50 of 13.4\% from DenseTeacher's fully supervised baseline (deng et al).


\subsection{The Paradigm Shift to Foundation Models and Zero-Shot Learning}
Looking beyond the iterative domain adaptation paradigm, the agricultural computer vision arena is witnessing a paradigm shift toward Visual Foundation Models (VFMs). After their introduction to the computer vision domain, vision transformers (ViTs) showed to have clear advantages over local CNNs in terms of capturing global relations \cite{dosovitskiy2020}. For instance, in a comparison between hierarchical ViTs (e.g. Swin Transformer) and standard ViT-B, the former achieved 92.3\% average accuracy on multi-crop settings \cite{ahmad2024}. As for novel architectures, the proposed WeedSwin model demonstrated outstanding performance for multi-stage classification tasks, attaining a test mAP of 0.993 at 218.27 FPS \cite{islam2025}. Meanwhile, to address the issue of large parameter budgets in most Transformers, a lightweight alternative was suggested, which slashed the number of parameters to 412K while managing to keep a reasonable inference time of 141ms for practical edge-level deployment \cite{yenna2026}.

Beyond the architecture design space, there has been an increasing number of works leveraging the powerful but simplistic Segment Anything Model (SAM) from Meta. For instance, the zero-shot adaptation ability of SAM was utilized for agricultural parcel boundary extraction in a phenology-aware manner \cite{simsek2026}. Furthermore, the same work used fine-tuned Stable Diffusion models to generate realistic synthetic imagery, thus minimizing the need for real image collection \cite{modak2024}. Due to their generic nature, foundation models lack the botanical knowledge to differentiate between crops and weeds; therefore, an interesting hybrid approach was introduced, which combined a lightweight object detector (YOLOv11) with SAM. Specifically, a few-shot fine-tuned YOLOv11 was utilized to generate bounding boxes, which, in turn, guided the segmentation performed by SAM2. This increased the mIOU by +11.3\% compared to a baseline while still maintaining a decent speed of 80.02 ms per image \cite{ghose2025}. In another hybrid SAM-YOLO approach, a Grounded YOLOv10-SAM2.1 model was used to segment cotton bolls, attaining an mIOU of 86.9\% at 18.9 ms per frame (approximately 52 FPS) (pooja et al).

Finally, while the VFMs demonstrated exceptional performance in many vision tasks, they show to have some shortcomings in the agriculture domain. For example, in a setting where the VFM is utilized as a zero-shot anomaly segmenter (i.e. consider weeds as anomalies), such models underperformed severely when facing extreme weed infestations. The reason for the failure was the skewed distribution of anomalies, which constituted the majority of the image in the most infested cases \cite{chong2025}. Lastly, a sanity check was done to see if foundational vision-language models such as Florence2 could be utilized for agriculture tasks without extra finetuning. However, the results showed that such an approach either underperformed significantly or performed worse than random (aashu et al).

\subsection{Multimodal Sensor Fusion and the Research Gap}
The repeated failures of both conventional CNNs and modern Foundation Models in highly occluded environments are dictated by the same limiting factor: the dependence on standard RGB cameras. Pure optical imagery is inherently limited in its ability to capture morphological details beyond the surface due to varying field illumination, shadows, and similar green hues of different plant species. An RGB-only segmentation model tested on the WeedsGalore dataset scored a dismal mIoU of 63.08\% (zeynep et al). Meanwhile, the fusion of different data streams poses its own challenges, including resolution mismatch, extreme temporal displacement between bands, and a lack of adaptive importance weighting (gao et al).

The way to overcome such limitations is to turn to alternative sources of information and attempt Multimodal sensor fusion. By combining the data from standard RGB, Near-Infrared (NIR), Red-Edge (RE), Thermal, or Short-Wave InfraRed (SWIR) sensors, one can obtain a more comprehensive understanding of the target scene. For instance, a simple fusion of RGB, NIR, and RE channels via a Transformer-CNN architecture allowed to boost the segmentation performance to 78.88\% mIUO with a mere 8.7 million parameters \cite{galymzhankyzy2025}. Another approach, implemented in the S3-Net framework, demonstrated the effectiveness of combining both RGB texture information and SWIR’s ability to peer through vegetation using Vision-Language Models. It achieved an incredible 96.9\% accuracy on seed phenotyping at 95 FPS (song et al). Despite the recent progress in self-supervised learning and contrastive vision models (MoCo v2, DenseCL), the majority of the studies show that self-supervised learning underperforms in comparison to conventional supervised training for plant phenotyping tasks. The reason for such a gap is tied to the ability of self-supervised models to capture spatial patterns, which is exacerbated by the high dimensionality and redundancy of plant imagery data \cite{ogidi2023}. To harness the power of such foundation models for plant phenotyping tasks, one has to implement aggressive model compression techniques. For instance, CDGLYOLOv11Lite utilizes operations fusion and TensorRT’s FP16/INT8 precision to achieve 67 FPS on a Jetson Xavier NX while only using 1.9 MB of memory \cite{mo2025}.

The analysis of the literature highlights the existence of a major Research Gap: while hybrid Foundation Models such as YOLOSAM address the issue of extensive manual annotation by enabling zero-shot segmentation, and Multimodal fusion helps to alleviate the ambiguity of optical data, the combination of the two has not been explored. In other words, there is little research on how to implement Multimodal fusion in a way that is computationally efficient enough for real-time edge processing. At the same time, most of the existing Foundation Models are optimized for unimodal RGB input. Thus, there is a need to develop a novel efficient multimodal foundation model that would be capable of zero-shot detection of various weed species in different crop fields.